% Part: history
% Chapter: biographies
% Section: rozsa-peter

\documentclass[../../../include/open-logic-section]{subfiles}

\begin{document}

\olfileid{his}{bio}{pet}

\olsection{روزا بيتر}

\olphoto{peter-rozsa}{R\'ozsa P\'eter}

وُلدت روزا بيتر، واسمها يومئذ روزا بوليتسر (R\'ozsa Politzer)، ببودابست
في بلاد المجر، في \OLClassicalDigits{17} فبراير \OLClassicalDigits{1905}. وإنما ذاع اسمها بما صنعت في الدوال
العودية، وكان عملها من الأصول التي قام عليها علم نظرية العودية.

ونشأت بيتر في عهد شديد الاضطراب---إذ كانت الحرب العالمية الأولى قائمة
وهي في حداثتها---ومع ذلك دخلت مدرسة ماريا تيريزيا الرفيعة للبنات في
بودابست، وفرغت منها سنة \OLClassicalDigits{1922}. ثم لحقت بجامعة باتسمايني بيتر في بودابست
(وهي التي سميت بعد ذلك جامعة إتفش لوراند، E\"otv\"os Lor\'and
University). وحملها إلحاح أبيها أولاً على درس الكيمياء، ثم عدلت إلى
الرياضيات وتخرجت سنة \OLClassicalDigits{1927}. وكانت مستوفية لأهلية تدريس الرياضيات في
المدارس الثانوية، غير أن الكساد الكبير أفسد اقتصاد العالم، فلم تجد عملاً
مستقراً؛ فكانت تقتات من أعمال متفرقة، تارة معلّمة خاصة وتارة مدرّسة
للرياضيات. ثم رجعت آخر الأمر إلى الجامعة لطلب الدراسات العليا في
الرياضيات. وكان عزمها أولاً أن تعمل في نظرية الأعداد، فلما علمت أن ما
انتهت إليه من النتائج قد سبق إلى برهانه غيرها أوشكت أن تدع الرياضيات
كلها. ثم حُثت على النظر في مبرهنتي غودل لعدم التمام، فأعادت البرهان على
طائفة من نتائجه بطرائق أخرى من غير أن تعلم بها. فعاد إليها وثوقها بنفسها،
وأخذت تكتب أول أبحاثها في نظرية العودية، مستضيئة ببرنامج ديفيد هيلبرت
التأسيسي. ونالت الدكتوراه سنة \OLClassicalDigits{1935}، ثم صارت سنة \OLClassicalDigits{1937} محررة في \emph{مجلة
المنطق الرمزي}.

وتُعد أبحاث بيتر الأولى من الأعمال المؤسسة لنظرية الدوال العودية؛ ففي
\cite{Peter1935a} فحصت الصلة بين أنواع شتى من العودية، ثم أثبتت في
\cite{Peter1935b} أن دالة معينة معرّفة بالعودية ليست من الدوال
العودية البدائية. وكان في ذلك تبسيط لنتيجة أقدم لفيلهلم أكرمان. وهذه
الدالة التي بسّطتها بيتر هي التي تسمى اليوم، في الغالب، دالة
أكرمان---والأدق أحياناً أن تسمى دالة أكرمان--بيتر. وهي صاحبة أول كتاب في
نظرية الدوال العودية \citep{Peter1951}.

وعلى جلالة عملها وسعة أثره، لم تُولَّ بيتر منصباً تعليمياً كاملاً إلا سنة
\OLClassicalDigits{1945}. وفي أثناء احتلال النازيين للمجر في الحرب العالمية الثانية منعتها
قوانينهم المعادية للسامية من التدريس. وأنشأت الحكومة سنة \OLClassicalDigits{1944} حياً معزولاً
ليهود بودابست، فصلته عن سائر المدينة وأحاطته بالحرس المسلح، فأُكرهت بيتر
على المقام فيه إلى أن حُرر سنة \OLClassicalDigits{1945}. ثم علّمت في كلية إعداد المعلمين
ببودابست، وانتقلت منذ سنة \OLClassicalDigits{1955} إلى جامعة إتفش لوراند. وكانت أول رياضية
مجرية تحرز درجة دكتور أكاديمي في الرياضيات، وأول امرأة تُنتخب في
الأكاديمية المجرية للعلوم.

وعُرفت بيتر معلّمة محبة للرياضيات، تؤثر أن تنقّب مع طلابها عن ماهية
المسائل الرياضية ووجوه حسنها على أن تلقي إليهم الكلام إلقاء؛ فكانوا
يسمونها مودةً «الخالة روزا». وتوفيت سنة \OLClassicalDigits{1977} وقد بلغت~\OLClassicalDigits{71} عاماً.

\begin{reading}
وللتوسع في سيرتها يُنظر \citep{Oconnor2014} و\citep{Andrasfai1986}، وفي
\citet{Tamassy1994} حديث موجز أجراه صاحبه مع بيتر. ومن رام قراءة ماتعة في
الرياضيات فليرجع إلى كتابها \emph{اللعب مع اللانهاية}
\citep{Peter2010}.
\end{reading}

\end{document}
